\documentclass[11pt]{article}
\usepackage{lmodern}
\usepackage[T1]{fontenc}
\usepackage{textcomp}
\usepackage[expansion=false]{microtype}
\usepackage{amsmath,amssymb}
\usepackage[version=4]{mhchem}
\usepackage{booktabs,longtable,array}
\usepackage{enumitem}
\usepackage{xcolor}
\usepackage{xurl}
\usepackage{hyperref}
\usepackage{geometry}
\geometry{margin=25mm}
\hypersetup{colorlinks=true,linkcolor=blue!50!black,urlcolor=blue!55!black,
  pdftitle={Hawthorn (Crataegus spp.): A Clinician-Oriented Review of the Past Year's Medical Literature (October 2025 - October 2026) with a Phytochemical Update},
  pdfauthor={Walker Kirkpatrick, ND},
  pdfkeywords={Crataegus, hawthorn, phytochemistry, procyanidins, cardiovascular, review},
  pdfsubject={Grounded literature review, October 2025 - October 2026}}
\newcommand{\src}[1]{[#1]}
\emergencystretch=3em
\setlength{\parindent}{0pt}
\setlength{\parskip}{0.55em}
\title{Hawthorn (\emph{Crataegus} spp.): A Clinician-Oriented Review of the Past Year's Medical Literature (October 2025--October 2026) with a Phytochemical Update}
\author{Walker Kirkpatrick, ND\thanks{Draft prepared with AI assistance (Hermes Agent); see notice below.}}
\date{October 25, 2026}
\begin{document}
\maketitle
\begin{center}\fbox{\parbox{0.93\textwidth}{\small Draft prepared with AI assistance (Hermes Agent) as literature-review support for the authoring clinician. Not medical advice; factual claims drawn from the literature carry numbered citations to retrieved sources, and the full evidence trail (verbatim quotes, ledger, fetched text) is preserved in the project repository. Clinical decisions remain the responsibility of the treating practitioner.}}\end{center}

\section*{Summary}
This review covers the peer-reviewed medical and pharmacological literature on hawthorn (\emph{Crataegus} spp.) published in the twelve months from October 25, 2025 through October 26, 2026, as indexed in PubMed/MEDLINE. A topic search returned 502 date-matched records, of which 89 carried \emph{Crataegus} or hawthorn in the title; after relevance triage, 52 focused reports ground this review, supplemented by PubChem records for compound structures and a ClinicalTrials.gov query for the trial pipeline.

The year's clinical yield was modest but real. A network meta-analysis of plant extracts in heart failure pooled 28 randomized trials (3,650 patients) and included hawthorn extract among 15 evaluated botanicals [1]. A randomized, double-blind, placebo-controlled crossover trial found that a \emph{Scutellaria baicalensis}/\emph{Crataegus laevigata} supplement with magnesium and chromium improved mood and stress-buffered cognition in subjectively stressed adults [4]. A single-blind randomized trial found hawthorn vinegar gargle reduced oral mucositis and improved oral health-related quality of life in people with diabetes [6].

The mechanistic yield was dominated by compound-level work. A systematic ex-vivo screen of 21 hawthorn constituents identified the methylated flavonols isorhamnetin and tamarixetin as the most active vasodilators, acting via L-type calcium-channel interference at concentrations well above typical human plasma levels [54]. Procyanidin C1 emerged as a senolytic-like anti-atherosclerotic whose direct target is HSP90 [15]; leucocyanidin as a dual AMPK/NRF2 activator against steatotic liver disease [69]; and vitexin as a multitarget antihyperuricemic acting on PPAR$\gamma$ and ABCG2 [16]. A new natural product, pinnatichromen, was described from \emph{C. pinnatifida} leaves [33].

Safety and quality signals warrant attention: a case report of acute diarrhea requiring emergency care after unsupervised use of a \emph{Crataegus mexicana} root-based weight-loss supplement [57]; a pharmacokinetic demonstration that hawthorn's ursolic acid alters transporter-mediated handling of a co-administered phytochemical [56]; and metabolomic confirmation that species, geography, and processing drive large composition differences that argue for chemotype-controlled products [68][31].

\section*{Scope and method}
Sources. PubMed was searched through NCBI E-utilities (esearch/esummary/efetch) with the term (\emph{Crataegus} OR hawthorn OR "WS 1442") restricted by publication date to 2025/10/25--2026/10/26, plus seven themed sub-queries spanning clinical trial publication types, phytochemistry, cardiovascular, metabolic, mechanistic, safety, and neuropsychiatric topics. Titles and abstracts were triaged for substantive focus on \emph{Crataegus}; 52 records were retained at abstract level. Chemical formulas and molecular weights in Table 1 were retrieved per compound from PubChem's PUG REST API, cited individually in the table.

This is a rapid, structured review, not a systematic review: one database, abstract-level grounding, no protocol registration, and AI-assisted triage. The trial pipeline statement is grounded in a ClinicalTrials.gov API v2 query of interventional studies naming \emph{Crataegus} or hawthorn [87]. Every claim carries a numbered citation to the retrieved record; verbatim supporting quotations for load-bearing claims are preserved in the project repository. The limitations of that scope are stated in the closing section.

\section*{Hallmark phytochemistry: the compounds of \emph{Crataegus}}
The year's reviews converge on a stable chemical picture in which flavonoids and oligomeric proanthocyanidins are the principal cardioactive constituents, exhibiting vasodilatory, antioxidant, mild positive inotropic, and antiarrhythmic properties [11]. A comprehensive UHPLC-Q-Exactive Orbitrap survey of 54 samples from four species, eleven production areas, and three drying methods identified 117 components spanning seven chemical classes, with phenolic acids and flavonoids the predominant constituents [68]. Reviews of the full composition list flavonoids, triterpenoids, organic acids, and polysaccharides as the major bioactive groups, alongside coverage of pharmacokinetics and adverse reactions [53][45].

Marker flavonoids. The flavone C-glycoside vitexin and the flavonol O-glycoside hyperoside remain the workhorse markers for species differentiation and quality control: principal-component analysis of eleven Sicilian \emph{Crataegus} taxa identified hyperoside and vitexin as the main contributors to sample differentiation, with hyperoside ranging up to 1.47 mg/g dry weight and total polyphenols peaking at 25.54 mg GAE/g in \emph{C. laciniata} [50]. Among three Turkish endemic species, \emph{C. peshmenii} leaf extract was richest in vitexin and hyperoside [13]. Free phenolics in \emph{C. azarolus} leaves reached 78.61 mg GAE/g dry weight, with a further 13.56 mg GAE/g bound fraction [22]; ethanol extracts of \emph{C. monogyna} carried 80.92 mg GAE/g total phenolics and 51.94 mg RE/g flavonoids [42]. GC-MS profiling of \emph{C. azarolus} flagged quercetin and the flavone C-glycoside lucenin-2 as constituents of proposed CNS relevance [59].

Procyanidins and flavanols. Hawthorn-leaf condensed tannins were confirmed by MALDI-TOF MS to consist mainly of procyanidins with smaller propelargonidin and prodelphinidin contributions [60]. Low-pressure chromatography of a procyanidin-rich \emph{C. monogyna} bark extract produced a fraction that was almost pure procyanidin B2 (88.8\% of total phenolics) and higher-polymer fractions whose antioxidant activity tracked procyanidin oligomer content [18]. A fraction containing 741.3 mg/g (-)-epicatechin anchored the low-molecular-weight end of that fractionation [18]. The first detailed phytochemical investigation of the endemic Kazakh \emph{C. almaatensis} isolated (-)-epicatechin, the flavolignan cinchonain Ia, kakispyrol, and two diaryl-ether glucosides, with cinchonain Ia showing radical-scavenging activity comparable to ascorbic acid [26].

Triterpenoids, polysaccharides, anthocyanins. Ursolic acid, hawthorn's signature triterpenoid acid, was shown to enhance intestinal absorption of a co-administered stilbene glycoside by downregulating efflux transporter MRP2 while upregulating hepatic UGT1A1 [56]. Maslinic acid in hawthorn juice was implicated as a driver of synergy with notoginseng stem-leaf saponins [63]. Polysaccharide chemistry advanced on two fronts: acidic, galacturonic-acid-rich leaf polysaccharides (10.19--435.52 kDa) whose most active fraction suppressed the SREBP-1c pathway in HepG2 cells [37], and habitat-level quality control showing that galacturonic-acid content and molecular weight (2.07--5.57 $\times$ 10$^6$ Da) correlate with radical-scavenging activity across ten growing regions [62]. The anthocyanin cyanidin 3-O-arabinoside was studied for metabolic effects, but note it derives from \emph{Malus doumeri}, commercially called "large-fruited hawthorn" and not a \emph{Crataegus} at all [64][34].

\begin{center}\small
\textbf{Table 1. Hallmark \emph{Crataegus} compounds with 2025--26 findings.} Formulas and molecular weights were retrieved per compound from PubChem; the rightmost column cites the year's finding and the matching PubChem record.\\[0.4em]
\setlength{\tabcolsep}{3pt}
\begin{longtable}{@{}>{\raggedright\arraybackslash}p{2.55cm}>{\raggedright\arraybackslash}p{2.5cm}>{\raggedright\arraybackslash}p{2.0cm}>{\raggedright\arraybackslash}p{1.15cm}>{\raggedright\arraybackslash}p{6.85cm}@{}}
\toprule
Compound & Class & Formula & MW & 2025--26 highlight \\
\midrule
\endhead
Vitexin & Flavone C-glycoside & \ce{C21H20O10} & 432.4 & Lead antihyperuricemic flavonoid; PPAR$\gamma$/ABCG2 [16][70] \\
Hyperoside & Flavonol O-glycoside & \ce{C21H20O12} & 464.4 & Chief species-differentiating marker [50][71] \\
Rutin & Flavonol O-glycoside & \ce{C27H30O16} & 610.5 & Cardioprotective food flavonoid panel [30][84] \\
Quercetin & Flavonol aglycone & \ce{C15H10O7} & 302.2 & CNS-relevant constituent of \emph{C. azarolus} [59][72] \\
Isorhamnetin & Methylated flavonol & \ce{C16H12O7} & 316.3 & Most active vasodilator in ex-vivo screen [54][73] \\
Tamarixetin & Methylated flavonol & \ce{C16H12O7} & 316.3 & Vasodilator; L-type Ca$^2$$^+$-channel blockade [54][74] \\
(-)-Epicatechin & Flavanol & \ce{C15H14O6} & 290.3 & 741.3 mg/g in bark fraction I [18][75] \\
(+)-Catechin & Flavanol & \ce{C15H14O6} & 290.3 & Predominant fruit phenolic after digestion [24][76] \\
Procyanidin B2 & Procyanidin dimer & \ce{C30H26O12} & 578.5 & 88.8\% of phenolics in bark fraction II [18][77] \\
Procyanidin C1 & Procyanidin trimer & \ce{C45H38O18} & 866.8 & Senolytic-like; direct HSP90 target [15][78] \\
Leucocyanidin & Flavan-3,4-diol & \ce{C15H14O7} & 306.3 & Dual AMPK/NRF2 activator in MASLD [69][79] \\
Ursolic acid & Triterpenoid acid & \ce{C30H48O3} & 456.7 & Alters MRP2/UGT1A1 handling of THSG [56][80] \\
Maslinic acid & Triterpenoid acid & \ce{C30H48O4} & 472.7 & Synergy driver with notoginseng saponins [63][81] \\
Chlorogenic acid & Phenolic acid ester & \ce{C16H18O9} & 354.3 & Major phenolic-acid class in 4-species survey [68][82] \\
Ellagic acid & Phenolic acid lactone & \ce{C14H6O8} & 302.2 & 109\% intestinal bioaccessibility (hawthorn) [32][83] \\
Cyanidin 3-O-arabinoside & Anthocyanin & \ce{C20H19O10+} & 419.4 & CD36/AMPK/PGC-1$\alpha$ axis (\emph{Malus doumeri}) [64][85] \\
Lucenin-2 & Flavone C-glycoside & \ce{C27H30O16} & 610.5 & GC-MS constituent with CNS hypotheses [59][86] \\
Pinnatichromen & Novel flavone glycoside & --- & --- & New natural product; inhibits IL-1$\beta$ production [33] \\
\bottomrule
\end{longtable}
\end{center}

\section*{New natural products and species-level discoveries}
A previously unreported flavone glycoside, pinnatichromen, was isolated from leaves of \emph{C. pinnatifida} var. \emph{major} and fully characterized by NMR and ESI-MS; it inhibited IL-1$\beta$ production at 20 $\mu$M [33]. The \emph{C. almaatensis} study above is the first detailed chemical report for that species [26], and first flavonoid profiles were published for the Turkish endemics \emph{C. peshmenii}, \emph{C. petrodavisii}, and \emph{C. christensenii}, with \emph{C. peshmenii} leaves richest in cardioactive markers [13]. Multi-technique characterization of \emph{C. monogyna} (LC-MS/MS plus DFT plus docking) quantified antioxidant and enzyme-inhibition profiles relevant to neurochemical and glycemic targets [42].

Biosynthesis genetics moved the field toward predictable chemotypes. A genome-wide analysis identified 215 MYB transcription factors in \emph{C. pinnatifida} and functionally validated CpMYB5 as a positive regulator of proanthocyanidin biosynthesis acting on the promoters of CpANR and CpLAR [65]. Chromosome-scale genomes of two cultivars (809.1 Mb and 760.5 Mb, scaffold N50 48.5 and 46.8 Mb) tied peel color to CpMYB114-CpANS anthocyanin regulation and seed hardness to CpMYB44-CpCOMT lignin biosynthesis [61].

Bioaccessibility work tempered raw-content claims. Under standardized INFOGEST simulated digestion, bioaccessible phenolics and antioxidant capacities fell substantially, though catechin remained the predominant bioaccessible phenolic across three \emph{Crataegus} species [24]. Flavanols proved the most digestion-stable class in \emph{C. monogyna} fruits, with epicatechin remaining 97--98\% bioaccessible and ellagic acid reaching approximately 109\% in the intestinal phase [32]. Freeze-drying preserved the highest antioxidant activity among processing methods [68].

\section*{Cardiovascular findings}
Heart failure. The year's most clinically consequential paper is a systematic review and network meta-analysis of plant extracts in heart failure, pooling 28 randomized trials and 3,650 patients across 15 botanicals, hawthorn extract among them [1]. On SUCRA ranking, water extract of \emph{Terminalia arjuna} bark ranked first for NYHA-class improvement (80.4\%), \emph{Salvia miltiorrhiza} first for LVEF improvement (89.7\%), oak wood extract first for quality of life (100\%), and \emph{Astragalus} first for the six-minute walk test (89.5\%) [1]. The published abstract does not report hawthorn-specific ranking statistics, so hawthorn's comparative position within this network cannot be stated from the retrievable record. The authors' overall conclusion is that plant extracts can serve as effective adjunctive therapies to improve cardiac function, quality of life, and prognosis, with cautious interpretation owing to limited studies per intervention and scarce head-to-head evidence [1].

Vasodilatory actives identified. A Fitoterapia screening study tested 21 hawthorn constituents and absorbable metabolites on porcine coronary arteries, rat aorta, and rat mesenteric arteries, motivated by the observation that hawthorn lowers arterial blood pressure in humans while the active component remains unidentified [54]. Isorhamnetin and tamarixetin were the most active vasodilators, with EC$_{50}$ values of 21.7 and 30.9 $\mu$M in porcine coronary arteries, 16.7 and 23.3 $\mu$M in rat aorta, and 22.2 and 16.8 $\mu$M in mesenteric arteries [54]. The relaxation was endothelium-independent, unaffected by K$^+$-channel or cGMP/PKG manipulation, and instead blocked L-type Ca$^2$$^+$-channel-driven vasoconstriction in a dose-dependent fashion [54]. The authors note frankly that these EC$_{50}$ values sit orders of magnitude above typical human plasma levels (tens to low hundreds of nM), so no single hawthorn compound yet explains the clinical blood-pressure effect at achievable exposures [54].

Atherosclerosis. Procyanidin C1 from hawthorn modulated the metabolic-inflammatory axis in ApoE-/- mice on a high-fat diet and in a D-galactose senescence model, improving serum lipids, reducing aortic root plaque, vascular senescence markers, and inflammatory adhesion molecules, suppressing arachidonic-acid metabolic enzymes, and --- via DARTS, LC-MS/MS, CETSA, and docking --- identifying HSP90 as a direct target while selectively clearing senescent cells [15]. Polydatin combined with hawthorn flavonoids reduced atherosclerotic lesion burden by 44.8\%, 49.29\%, and 72.99\% at low, medium, and high doses in ApoE-/- mice, lowering total cholesterol, triglycerides, LDL, and VLDL, and acting through a YTHDF2-dependent suppression of the mitochondrial fission effector DRP1 [25]. A seven-herb food-medicine formulation containing \emph{C. pinnatifida} (Yaoshi Tongyuan Tablet) attenuated dyslipidemia and inflammation via PI3K/Akt/NF-$\kappa$B signaling, with kaempferol, isorhamnetin, and quercetin identified as central actives on AKT1 [29].

Reviews and tissue-level support. Comprehensive reviews published this year restate the mechanistic case --- antioxidant, anti-inflammatory, hypolipidemic, antihyperglycemic, antihypertensive, anticoagulant, anti-ischemic, diuretic, and cardiotonic --- for \emph{Crataegus} species [11][12][51]. The geroprotection frame is new and notable: \emph{Crataegus} phytochemicals are proposed to act on oxidative stress, inflammaging, mitochondrial dysfunction, cellular senescence, and the senescence-associated secretory phenotype [51]. At tissue level, leaf extract of \emph{C. peshmenii} reduced inflammatory and oxidative markers released from isolated rat thoracic aorta [13], and in a veterinary high-altitude model, \emph{C. oxyacantha} extract in drinking water raised nitric oxide by 57.1\%, SOD by 16.4\%, and GPx by 20.6\% while lowering hematocrit 13.1\% and improving growth efficiency in cold-stressed broilers [47]. Food-medicine-continuum reviews position hawthorn fruit among the leading anti-atherosclerotic materia medica [35].

\section*{Metabolic, hepatic, and renal findings}
Steatotic liver disease was the year's busiest front. Leucocyanidin, a hawthorn flavonoid, emerged as a dual-target agent in a 60\%-fat-diet mouse model and a free-fatty-acid HepG2 model: dose-dependently reducing liver weight, serum triglycerides and cholesterol, and IL-6/TNF-$\alpha$, restoring fatty-acid $\beta$-oxidation through the NRF2 antioxidant axis, and suppressing lipogenesis and fatty-acid uptake through direct allosteric binding to AMPK [69]. Hawthorn procyanidin extract attenuated NAFLD in mice by remodeling bile-acid profiles ($\omega$MCA, T$\alpha$MCA, T$\beta$MCA, and DCA) through the gut microbiota, activating hepatic FXR and CYP7A1, and downregulating SREBP1 via the FXR-FGF15 enterohepatic pathway [17]. A 70\% ethanol fruit extract characterized as 17 primary compounds inhibited ERK and regulated hepatic purine metabolism in NAFLD mice and HepG2 cells [55]. Reviews synthesize the mechanism set --- suppressed hepatic lipogenesis, enhanced fatty-acid $\beta$-oxidation, improved insulin signaling, adipose and bile-acid regulation, and gut-microbiota reshaping --- and conclude the evidence base remains preclinical with limited clinical investigation [45][40].

Polysaccharides carried much of the hepatic story. North hawthorn polysaccharide (200--800 mg/kg/day for 12 weeks) dose-dependently alleviated weight gain, dyslipidemia, steatosis, and inflammation, activated the Nrf2/Keap1/HO-1 axis, and restored \emph{Lactobacillus}, \emph{Bifidobacterium}, and \emph{Akkermansia muciniphila} abundance in mice [46]. Hawthorn polysaccharides lowered the Firmicutes/Bacteroidetes ratio, raised HDL-C, and lowered LDL-C in high-fat-fed mice [67], and leaf polysaccharide HLP-1 lowered hepatocyte triglycerides and cholesterol while suppressing SREBP-1c [37].

Urate. Vitexin was identified from 17 hawthorn flavonoid structures (UPLC-MS/MS) by SwissADME prioritization as the lead antihyperuricemic: in hyperuricemic mice it reduced serum uric acid by inhibiting hepatic xanthine oxidase (XOD) and ADA, upregulated renal and intestinal ABCG2 toward normal, modulated OAT1 and GLUT9, and protected liver, kidney, and intestine while restoring epithelial barrier integrity, consistent with its predicted PPAR$\gamma$/ABCG2 mechanism [16].

Glycemic targets. Hawthorn-leaf condensed tannins inhibited $\alpha$-glucosidase with an IC$_{50}$ of 8.0 $\mu$g/mL (among the stronger botanical values reported this year) through noncovalent binding that perturbs the tyrosine/tryptophan microenvironment and enzyme secondary structure, and also suppressed nonenzymatic glycation products, protein carbonylation, and amyloid cross-$\beta$ structures [60]. $\alpha$-Glucosidase and $\alpha$-amylase inhibition was likewise demonstrated for \emph{C. azarolus} leaf phenolics [22], \emph{C. monogyna} extracts [42], and machine-learning-optimized processed hawthorn vinegar [48].

Gastrointestinal inflammation. \emph{C. monogyna} fruit extract (100--400 mg/kg, 10 days) dose-dependently protected colonic mucosa in an acetic-acid colitis model, restoring MDA, thiol, SOD, CAT, and GPx toward control and normalizing CRP and ALP [66]. Compound-level gut-axis pharmacology extended beyond the liver: hyperoside, a flavonol glycoside abundant in hawthorn fruits, protected against \emph{E. coli} endometritis in mice through gut-microbiota-dependent production of hydroxyphenyllactic acid acting on TLR4 [28].

Combinations and caveats. In ovariectomized rats, combining high-intensity interval training with \emph{C. pentagyna} fruit extract reduced visceral fat by 72.8\%, versus 33.7\% for training alone and 65.5\% for extract alone, without dose-dependence across 100--300 mg/kg [58]. Hawthorn juice plus notoginseng stem-leaf saponins improved glucose and lipid parameters and blood viscosity in a high-fat, high-sugar model, with maslinic acid implicated in saponin transformation and synergy [63]. Readers should note that the cyanidin 3-O-arabinoside metabolic study, however favorable, concerns \emph{Malus doumeri}, a species marketed as "large-fruited hawthorn" and taxonomically distinct from \emph{Crataegus} [64][34].

\section*{Neuropsychiatric and neuroprotective findings}
Stress and cognition RCT. In a randomized, double-blind, placebo-controlled crossover trial (43 enrolled, 35 analyzed; 15 days per arm; NCT05757050), a supplement combining \emph{S. baicalensis} root and \emph{C. laevigata} fruit extracts with magnesium and chromium acutely reduced serial-3s subtraction errors, improved serial-7s working-memory performance under a multitasking stressor, and chronically reduced total mood disturbance, anger/hostility, and trait anxiety, with no serious adverse effects [4]. Attribution to \emph{Crataegus} specifically is not possible from this four-ingredient design --- the trial supports the combination product, not any single botanical [4].

Neuroprotective mechanisms. Total flavonoids of hawthorn leaves (TFHL), a traditional Chinese medicine preparation, promoted motor recovery and preserved neurons after spinal cord injury in rats while inhibiting TLR4 complex formation, NF-$\kappa$B nuclear translocation, and IL-1$\beta$, IL-6, and TNF-$\alpha$ mRNA expression in microglia [49]. A companion network-pharmacology study mapped 15 active TFHL compounds to 69 targets enriched in inflammation, oxidative stress, and apoptosis pathways, with MAPK signaling validated in BV-2 microglia and RAW264.7 macrophages [36]. These are early preclinical signals in a new indication area for hawthorn leaf flavonoids [49][36].

Psychiatric hypotheses and a methodological caution. A chromatographic study of \emph{C. azarolus} proposed psychiatric implications from its chemistry, highlighting quercetin and lucenin-2 as constituents with preclinical anxiolytic and antidepressant-like signals [59]. The same GC-MS analysis tentatively identified "digitoxin" among major compounds --- an implausible finding for a rosaceous fruit that illustrates why library-match GC-MS identifications require orthogonal confirmation before any clinical inference [59].

\section*{Safety, quality, and interactions}
Adverse-event signal. A Colombian case report describes a 41-year-old woman who presented to the emergency department with 15 days of daily diarrhea, epigastric pain, and moderate dehydration after unsupervised use of a \emph{Crataegus mexicana} root-based weight-loss supplement; she recovered with intravenous fluids and analgesics within 48 hours [57]. The report is a useful reminder that "natural" weight-loss products --- here a root-based preparation, not the well-studied leaf and flower extracts --- carry real risk when used without supervision [57].

Pharmacokinetic interactions. Hawthorn's ursolic acid increased Cmax of the \emph{Polygonum multiflorum} constituent THSG while decreasing Tmax and terminal half-life, by downregulating intestinal efflux transporter MRP2 and upregulating hepatic UGT1A1, with reduced serum exosomal miR-150-5p [56]. This is direct evidence that a hawthorn triterpenoid can alter the disposition of a co-administered phytochemical --- mechanistically relevant to the multi-botanical stacking common in practice [56]. A dedicated digoxin--hawthorn pharmacokinetic interaction trial (NCT00006330) is registered on ClinicalTrials.gov as completed, part of a registry record of 37 interventional studies matching \emph{Crataegus} or hawthorn [87].

Quality and variability. Species, production area, and drying method each shifted the metabolome across 117 identified components, and vacuum freeze-drying preserved the highest antioxidant activity [68]. Polysaccharide chemistry and bioactivity varied across ten habitats with structure--activity correlations [62]; marker content varied several-fold across eleven Sicilian taxa [50]; extraction solvent shifted both yield and bioactivity in \emph{C. monogyna} [42] and \emph{C. azarolus} [22]. Reviews this year uniformly call for standardization: chemotype-controlled, preparation-defined formulations with PK/PD-guided dosing [31], standardization of bioactive components before efficacy and optimal-dose claims [41], and quality frameworks spanning germplasm to processed product [19]. Methodological guidance for flavonoid determination in cardioprotective foods, including hawthorn, was reviewed in detail [30].

Nomenclature and authentication. "Hawthorn" in commerce does not guarantee \emph{Crataegus}: the metabolomics and cyanidin 3-O-arabinoside papers this year concern \emph{Malus doumeri}, marketed as large-fruited hawthorn [34][64]. Combined with tentative GC-MS identifications [59] and species-dependent content [68][50], the year's literature makes botanical authentication and chemotype declaration a patient-safety issue, not merely a research nicety [31].

\section*{Trial pipeline}
A ClinicalTrials.gov API query for interventional studies naming \emph{Crataegus} or hawthorn returned 37 registrations: 23 completed, 9 of unknown status, 2 terminated, 2 not yet recruiting, and 1 enrolling by invitation [87]. The single active dedicated trial is NCT07166965, ketone plus hawthorn extract supplementation in congestive heart failure, enrolling by invitation [87]. Two not-yet-recruiting polyherbal protocols include hawthorn-containing formulas, including Longqi Jiangzhi decoction for NAFLD (NCT07320456) [87]. The registry anchors for hawthorn cardiology --- the HERB CHF chronic heart failure trial (NCT00343902), the digoxin interaction study (NCT00006330), and a nitric-oxide-mediated vasodilation study of standardized extract (NCT01331486) --- are all completed [87]. In short, the registered pipeline holds no large new dedicated hawthorn trial, and near-term evidence generation appears to be riding on small combination-product studies [87].

\section*{Clinical synthesis for practice}
The paragraphs below are the authoring clinician's synthesis, labeled as interpretation rather than cited fact.

Heart failure. The new network meta-analysis keeps plant extracts, hawthorn included, in the adjunctive conversation, but its headline rankings favor other botanicals, and no new dedicated hawthorn heart-failure trial appeared in this window [1][87]. Where hawthorn is used, the evidence base remains the historical leaf and flower extract literature; nothing this year changes that (interpretation).

Blood pressure. The vasodilator screening sharpens the target --- methylated flavonols acting on vascular smooth-muscle calcium channels --- while demonstrating that no single compound reaches active plasma concentrations, which argues either for whole-extract synergy or for formulation work rather than single-ingredient supplementation [54] (interpretation).

Metabolic and hepatic indications. The steatotic-liver mechanistic base (procyanidins, leucocyanidin, polysaccharide gut--liver signaling) grew substantially this year, but every relevant study was in mice or cells [69][17][15][46]; no human metabolic endpoint trial appeared in this window, so hawthorn's metabolic role remains adjunctive and unproven in humans (interpretation).

Stress, mood, and topical use. One positive multi-ingredient RCT supports a \emph{Crataegus laevigata}-containing combination for stress and mood [4]; one positive RCT supports hawthorn vinegar gargle for diabetic oral mucositis [6]; neither licenses monotherapy claims for \emph{Crataegus} (interpretation).

Product selection. Given demonstrated species-, geography-, and processing-driven composition differences, clinicians should prefer chemotype-declared, standardized extracts and treat authenticity --- \emph{Crataegus} versus marketed "hawthorns" from other genera --- as a first-order selection criterion [68][31][34] (interpretation).

Safety. This year added a root-product case report [57] and a transporter-interaction mechanism [56] to the long-standing benign-tolerability record [4]; the practical takeaways are to avoid unregulated root-based weight-loss products, to review concomitant botanical stacks, and to remember the historical digoxin co-administration question when hawthorn is used alongside narrow-therapeutic-index cardiotonics [87] (interpretation).

\section*{Limitations}
This review is deliberately scoped and its limits are explicit. It covers PubMed/MEDLINE only, so literature in databases or languages outside that index is not represented; grounding is at abstract level, with full texts not retrieved; triage was AI-assisted; the base literature is dominated by preclinical and food-science designs with heterogeneous species, plant parts, extracts, and doses, limiting synthesis [1]; several high-visibility items report hawthorn-inclusive rather than hawthorn-specific statistics [1]; and in-vitro potency values frequently exceed achievable human exposures [54]. Counts (502 records matched; 89 hawthorn-titled; 52 retained) are method outputs from this search, not bibliometric claims. No load-bearing claim in this review required an unverified marker; where the retrievable record could not support a claim, the claim was not made.

\section*{Sources}
\footnotesize Entries are numbered by retrieval identity, matching the inline [n] markers inline. Every cited entry below carries a captured verbatim supporting quotation in the project repository ledger (citations/ledger.json), alongside the fetched page text (research/pages/). PubMed abstracts are the content of record; two retained items lacking usable abstracts were not cited.
{\footnotesize
\noindent\hangindent=2.6em\hangafter=1 [1] Efficacy of plant extracts in heart failure patients: a systematic review and network meta-analysis.. \url{https://pubmed.ncbi.nlm.nih.gov/42277690}\par
\noindent\hangindent=2.6em\hangafter=1 [4] Effects of a Scutellaria baicalensis/Crataegus laevigata, magnesium and chromium supplement on stressed individuals: A r. \url{https://pubmed.ncbi.nlm.nih.gov/41194549}\par
\noindent\hangindent=2.6em\hangafter=1 [6] The effect of hawthorn vinegar, black mulberry syrup, and green tea on oral mucositis and oral health-related quality of. \url{https://pubmed.ncbi.nlm.nih.gov/42694165}\par
\noindent\hangindent=2.6em\hangafter=1 [11] Hawthorn (Crataegus Species): Phytochemical Insights and Cardiotonic Mechanisms in Cardiovascular Health- A Comprehensiv. \url{https://pubmed.ncbi.nlm.nih.gov/42473213}\par
\noindent\hangindent=2.6em\hangafter=1 [12] A Comprehensive Review on Cardioprotective Biological Properties of Crataegus orientalis in Post-COVID-19 Pandemic Cardi. \url{https://pubmed.ncbi.nlm.nih.gov/42661384}\par
\noindent\hangindent=2.6em\hangafter=1 [13] Innovative findings on three endemic Crataegus spp. from T{\"u}rkiye: flavonoid-enriched extracts with cardiovascular benefi. \url{https://pubmed.ncbi.nlm.nih.gov/40232171}\par
\noindent\hangindent=2.6em\hangafter=1 [15] Procyanidin C1 from hawthorn modulates metabolic-inflammatory axis to combat vascular senescence in atherosclerosis.. \url{https://pubmed.ncbi.nlm.nih.gov/41652755}\par
\noindent\hangindent=2.6em\hangafter=1 [16] Antihyperuricemic Effect of Hawthorn Flavonoid Vitexin: LC-MS Analysis, Network Pharmacology, and In Vivo Verification o. \url{https://pubmed.ncbi.nlm.nih.gov/42255677}\par
\noindent\hangindent=2.6em\hangafter=1 [17] The hawthorn (Crataegus pinnatifida) procyanidin extract attenuates nonalcoholic fatty liver disease in mice via remodel. \url{https://pubmed.ncbi.nlm.nih.gov/42300918}\par
\noindent\hangindent=2.6em\hangafter=1 [18] Phenolic Profile and Antioxidant Activity of Fractions of Procyanidin-Rich Hawthorn (Crataegus monogyna Jacq.) Bark Extr. \url{https://pubmed.ncbi.nlm.nih.gov/41302433}\par
\noindent\hangindent=2.6em\hangafter=1 [19] The Hawthorn (Crataegus spp.) Value Chain: An Integrated Analysis of Resource Availability, Phytochemical Characterizati. \url{https://pubmed.ncbi.nlm.nih.gov/41517214}\par
\noindent\hangindent=2.6em\hangafter=1 [22] Ultrasound-Assisted Extraction of Free and Bound Phenolics from Hawthorn (Crataegus azarolus) Leaves and their Antioxida. \url{https://pubmed.ncbi.nlm.nih.gov/42141147}\par
\noindent\hangindent=2.6em\hangafter=1 [24] Bioaccessibility of Phytochemicals, Amino Acids, Antioxidant Capacity, and Enzyme Inhibitory Activities of Hawthorn (Cra. \url{https://pubmed.ncbi.nlm.nih.gov/42559186}\par
\noindent\hangindent=2.6em\hangafter=1 [25] Polydatin in combination with hawthorn flavonoids alleviates atherosclerosis by inhibiting mitochondrial division throug. \url{https://pubmed.ncbi.nlm.nih.gov/42595105}\par
\noindent\hangindent=2.6em\hangafter=1 [26] Phytochemical Investigation of Crataegus almaatensis Leaves: Isolation, Structural Characterization, and Antioxidant Act. \url{https://pubmed.ncbi.nlm.nih.gov/42792188}\par
\noindent\hangindent=2.6em\hangafter=1 [28] Hyperoside, a dietary flavonoid, protects against endometritis via gut microbiota-dependent production of hydroxyphenyll. \url{https://pubmed.ncbi.nlm.nih.gov/41369516}\par
\noindent\hangindent=2.6em\hangafter=1 [29] A food-medicine homology formulation ameliorates atherosclerosis by attenuating dyslipidemia and inflammation via the PI. \url{https://pubmed.ncbi.nlm.nih.gov/41424781}\par
\noindent\hangindent=2.6em\hangafter=1 [30] [Methodological aspects of determining flavonoids in food products with cardioprotective potential -- a review of HPLC co. \url{https://pubmed.ncbi.nlm.nih.gov/41504733}\par
\noindent\hangindent=2.6em\hangafter=1 [31] Medicinal Plants and Phytochemicals in Cardioprotection-Mechanistic Pathways and Translational Roadmap.. \url{https://pubmed.ncbi.nlm.nih.gov/41598329}\par
\noindent\hangindent=2.6em\hangafter=1 [32] Bromatological Profile of Fruits from Sorbus aucuparia and Crataegus monogyna: Polyphenol Bioaccessibility and Inhibitor. \url{https://pubmed.ncbi.nlm.nih.gov/41897495}\par
\noindent\hangindent=2.6em\hangafter=1 [33] A novel flavonoid from leaves of Crataegus pinnatifida Bge. var. major and its anti-inflammatory activity.. \url{https://pubmed.ncbi.nlm.nih.gov/42054581}\par
\noindent\hangindent=2.6em\hangafter=1 [34] Comparative Metabolomic and Transcriptomic Analyses Identify Candidate Genes Associated with Flavonoid Accumulation and. \url{https://pubmed.ncbi.nlm.nih.gov/42280163}\par
\noindent\hangindent=2.6em\hangafter=1 [35] Bioactive Constituents, Mechanisms, and Complementary Therapeutic Applications of Food-Medicine Continuum Materia Medica. \url{https://pubmed.ncbi.nlm.nih.gov/42356474}\par
\noindent\hangindent=2.6em\hangafter=1 [36] Molecular mechanism of total flavonoids of hawthorn leaves in spinal cord injury treatment: A combined network pharmacol. \url{https://pubmed.ncbi.nlm.nih.gov/42894930}\par
\noindent\hangindent=2.6em\hangafter=1 [37] Structural characterization, in vitro antioxidant and lipid-lowering activities of purified polysaccharides from hawthor. \url{https://pubmed.ncbi.nlm.nih.gov/41052502}\par
\noindent\hangindent=2.6em\hangafter=1 [40] Recent advances in hawthorn (Crataegus L.) polysaccharides for prevention and treatment of non-alcoholic fatty liver dis. \url{https://pubmed.ncbi.nlm.nih.gov/41338509}\par
\noindent\hangindent=2.6em\hangafter=1 [41] Hawthorn (Crataegus monogyna Jacq.): A Review of Therapeutic Potential and Applications.. \url{https://pubmed.ncbi.nlm.nih.gov/41599280}\par
\noindent\hangindent=2.6em\hangafter=1 [42] Phytochemical Composition and Multi-Target Biological Activities of Crataegus monogyna: LC-MS/MS, DFT, and Molecular Doc. \url{https://pubmed.ncbi.nlm.nih.gov/41643173}\par
\noindent\hangindent=2.6em\hangafter=1 [45] Regulatory effects of hawthorn on lipid metabolic homeostasis: mechanisms, evidences, and perspectives.. \url{https://pubmed.ncbi.nlm.nih.gov/41889722}\par
\noindent\hangindent=2.6em\hangafter=1 [46] Gut-liver axis mediates the ameliorative effects of North Hawthorn polysaccharide on liver inflammation and oxidative st. \url{https://pubmed.ncbi.nlm.nih.gov/42077781}\par
\noindent\hangindent=2.6em\hangafter=1 [47] Crataegus oxyacantha Extract Enhances Cardiac Adaptation and Reduces Oxidative-Inflammatory Stress in Broiler Chickens R. \url{https://pubmed.ncbi.nlm.nih.gov/42295739}\par
\noindent\hangindent=2.6em\hangafter=1 [48] Machine learning-assisted optimization of ultrasound-ohmic processing of hawthorn vinegar: antidiabetic activity, phenol. \url{https://pubmed.ncbi.nlm.nih.gov/42488216}\par
\noindent\hangindent=2.6em\hangafter=1 [49] Total Flavonoids from Hawthorn Leaves Alleviate Acute Inflammatory Responses Following Spinal Cord Injury by Downregulat. \url{https://pubmed.ncbi.nlm.nih.gov/42634323}\par
\noindent\hangindent=2.6em\hangafter=1 [50] Polyphenolic Characterization, Antioxidant and Antihyperglycemic Activities of Crataegus spp. Fruit Extracts.. \url{https://pubmed.ncbi.nlm.nih.gov/42739450}\par
\noindent\hangindent=2.6em\hangafter=1 [51] Crataegus Species as a Source of Geroprotective Phytochemicals: Biological Activities, Molecular Mechanisms and Therapeu. \url{https://pubmed.ncbi.nlm.nih.gov/42795343}\par
\noindent\hangindent=2.6em\hangafter=1 [53] Research Progress on The Chemical Composition And Pharmacological Effects of Hawthron (Crataegus spp.): A Review.. \url{https://pubmed.ncbi.nlm.nih.gov/41431296}\par
\noindent\hangindent=2.6em\hangafter=1 [54] Detection of active vasodilatory compounds from hawthorn: Screening and determination of mechanism of action.. \url{https://pubmed.ncbi.nlm.nih.gov/41794075}\par
\noindent\hangindent=2.6em\hangafter=1 [55] Hawthorn fruit extracts alleviate high-fat diet-induced NAFLD by inhibiting ERK and regulating hepatic purine metabolism. \url{https://pubmed.ncbi.nlm.nih.gov/41744261}\par
\noindent\hangindent=2.6em\hangafter=1 [56] Enhancement of ursolic acid (UA), an active ingredient in hawthorn, on the uptake and metabolism of 2,3,5,4'-tetrahydrox. \url{https://pubmed.ncbi.nlm.nih.gov/42215098}\par
\noindent\hangindent=2.6em\hangafter=1 [57] Crataegus mexicana root-based supplement intoxication in a woman with weight loss persistent preoccupation: Case report.. \url{https://pubmed.ncbi.nlm.nih.gov/41875453}\par
\noindent\hangindent=2.6em\hangafter=1 [58] High-intensity interval training and Crataegus pentagyna extract enhance metabolic and oxidant profiles in ovariectomize. \url{https://pubmed.ncbi.nlm.nih.gov/42163902}\par
\noindent\hangindent=2.6em\hangafter=1 [59] Chromatographic analysis of Crataegus azarolus and possible psychiatric implications.. \url{https://pubmed.ncbi.nlm.nih.gov/41855569}\par
\noindent\hangindent=2.6em\hangafter=1 [60] Inhibition of Hawthorn Leaf Condensed Tannins on $\alpha$-Glucosidase and Nonenzymatic Glycosylation: Molecular Mechanisms of A. \url{https://pubmed.ncbi.nlm.nih.gov/41569292}\par
\noindent\hangindent=2.6em\hangafter=1 [61] Integrated genomic and transcriptomic analyses reveal the genetic and molecular mechanisms underlying hawthorn peel colo. \url{https://pubmed.ncbi.nlm.nih.gov/40220858}\par
\noindent\hangindent=2.6em\hangafter=1 [62] Quality control and spectrum-effect relationship of Crataegus pinnatifida polysaccharide from different habitats.. \url{https://pubmed.ncbi.nlm.nih.gov/40078042}\par
\noindent\hangindent=2.6em\hangafter=1 [63] Hawthorn juice transforms notoginseng stem-leaf saponins: Synergistically regulates blood viscosity and ameliorates glyc. \url{https://pubmed.ncbi.nlm.nih.gov/41338313}\par
\noindent\hangindent=2.6em\hangafter=1 [64] Cyanidin-3-O-arabinoside improves high-fat diet-induced glucose and lipid metabolism disorder in mice by restoring fatty. \url{https://pubmed.ncbi.nlm.nih.gov/42305060}\par
\noindent\hangindent=2.6em\hangafter=1 [65] Genome-wide identification of the MYB family in Crataegus pinnatifida and functional analysis of CpMYB5 involved in proa. \url{https://pubmed.ncbi.nlm.nih.gov/41077282}\par
\noindent\hangindent=2.6em\hangafter=1 [66] Preventive Effects of Crataegus monogyna Fruit Extract on Oxidative Damage and Inflammatory Response in a Rat Model of A. \url{https://pubmed.ncbi.nlm.nih.gov/42440624}\par
\noindent\hangindent=2.6em\hangafter=1 [67] Hawthorn Polysaccharides Relieve High-Fat Blood Disease in Mice by Improving Intestinal Flora and Promoting Lipid Metabo. \url{https://pubmed.ncbi.nlm.nih.gov/42121468}\par
\noindent\hangindent=2.6em\hangafter=1 [68] Effects of species, production area and processing methods on chemical composition and antioxidant activity of hawthorn.. \url{https://pubmed.ncbi.nlm.nih.gov/42661862}\par
\noindent\hangindent=2.6em\hangafter=1 [69] Novel dual AMPK/NRF2 activation by leucocyanidin from Hawthorn (Crataegus) for mitochondria repair-Targeted therapy of h. \url{https://pubmed.ncbi.nlm.nih.gov/41365192}\par
\noindent\hangindent=2.6em\hangafter=1 [70] PubChem: vitexin. \url{https://pubchem.ncbi.nlm.nih.gov/rest/pug/compound/name/vitexin/property/MolecularFormula,MolecularWeight/JSON}\par
\noindent\hangindent=2.6em\hangafter=1 [71] PubChem: hyperoside. \url{https://pubchem.ncbi.nlm.nih.gov/rest/pug/compound/name/hyperoside/property/MolecularFormula,MolecularWeight/JSON}\par
\noindent\hangindent=2.6em\hangafter=1 [72] PubChem: quercetin. \url{https://pubchem.ncbi.nlm.nih.gov/rest/pug/compound/name/quercetin/property/MolecularFormula,MolecularWeight/JSON}\par
\noindent\hangindent=2.6em\hangafter=1 [73] PubChem: isorhamnetin. \url{https://pubchem.ncbi.nlm.nih.gov/rest/pug/compound/name/isorhamnetin/property/MolecularFormula,MolecularWeight/JSON}\par
\noindent\hangindent=2.6em\hangafter=1 [74] PubChem: tamarixetin. \url{https://pubchem.ncbi.nlm.nih.gov/rest/pug/compound/name/tamarixetin/property/MolecularFormula,MolecularWeight/JSON}\par
\noindent\hangindent=2.6em\hangafter=1 [75] PubChem: epicatechin. \url{https://pubchem.ncbi.nlm.nih.gov/rest/pug/compound/name/epicatechin/property/MolecularFormula,MolecularWeight/JSON}\par
\noindent\hangindent=2.6em\hangafter=1 [76] PubChem: catechin. \url{https://pubchem.ncbi.nlm.nih.gov/rest/pug/compound/name/catechin/property/MolecularFormula,MolecularWeight/JSON}\par
\noindent\hangindent=2.6em\hangafter=1 [77] PubChem: procyanidin B2. \url{https://pubchem.ncbi.nlm.nih.gov/rest/pug/compound/name/procyanidin%20B2/property/MolecularFormula,MolecularWeight/JSON}\par
\noindent\hangindent=2.6em\hangafter=1 [78] PubChem: procyanidin C1. \url{https://pubchem.ncbi.nlm.nih.gov/rest/pug/compound/name/procyanidin%20C1/property/MolecularFormula,MolecularWeight/JSON}\par
\noindent\hangindent=2.6em\hangafter=1 [79] PubChem: leucocyanidin. \url{https://pubchem.ncbi.nlm.nih.gov/rest/pug/compound/name/leucocyanidin/property/MolecularFormula,MolecularWeight/JSON}\par
\noindent\hangindent=2.6em\hangafter=1 [80] PubChem: ursolic acid. \url{https://pubchem.ncbi.nlm.nih.gov/rest/pug/compound/name/ursolic%20acid/property/MolecularFormula,MolecularWeight/JSON}\par
\noindent\hangindent=2.6em\hangafter=1 [81] PubChem: maslinic acid. \url{https://pubchem.ncbi.nlm.nih.gov/rest/pug/compound/name/maslinic%20acid/property/MolecularFormula,MolecularWeight/JSON}\par
\noindent\hangindent=2.6em\hangafter=1 [82] PubChem: chlorogenic acid. \url{https://pubchem.ncbi.nlm.nih.gov/rest/pug/compound/name/chlorogenic%20acid/property/MolecularFormula,MolecularWeight/JSON}\par
\noindent\hangindent=2.6em\hangafter=1 [83] PubChem: ellagic acid. \url{https://pubchem.ncbi.nlm.nih.gov/rest/pug/compound/name/ellagic%20acid/property/MolecularFormula,MolecularWeight/JSON}\par
\noindent\hangindent=2.6em\hangafter=1 [84] PubChem: rutin. \url{https://pubchem.ncbi.nlm.nih.gov/rest/pug/compound/name/rutin/property/MolecularFormula,MolecularWeight/JSON}\par
\noindent\hangindent=2.6em\hangafter=1 [85] PubChem: Cyanidin 3-O-arabinoside. \url{https://pubchem.ncbi.nlm.nih.gov/rest/pug/compound/name/Cyanidin%203-O-arabinoside/property/MolecularFormula,MolecularWeight/JSON}\par
\noindent\hangindent=2.6em\hangafter=1 [86] PubChem: lucenin-2. \url{https://pubchem.ncbi.nlm.nih.gov/rest/pug/compound/name/lucenin-2/property/MolecularFormula,MolecularWeight/JSON}\par
\noindent\hangindent=2.6em\hangafter=1 [87] ClinicalTrials.gov API v2: interventional trials, Crataegus OR hawthorn. \url{https://clinicaltrials.gov/api/v2/studies?query.intr=%28Crataegus%20OR%20hawthorn%29&fields=NCTId%2CBriefTitle%2COverallStatus%2CStudyType%2CStartDate%2CEnrollmentCount&pageSize=100}\par
\vspace{0.3em}\noindent Retrieved October 25, 2026 via NCBI E-utilities, PubChem PUG REST, and ClinicalTrials.gov API v2.}
\end{document}
